
\documentclass{article}
\usepackage{colortbl}
\usepackage{makecell}
\usepackage{multirow}
\usepackage{supertabular}

\begin{document}

\newcounter{utterance}

\centering \large Interaction Transcript for game `cladder', experiment `full\_v1.5\_default', episode 2646 with qwen3.5{-}27b.
\vspace{24pt}

{ \footnotesize  \setcounter{utterance}{1}
\setlength{\tabcolsep}{0pt}
\begin{supertabular}{c@{$\;$}|p{.15\linewidth}@{}p{.15\linewidth}p{.15\linewidth}p{.15\linewidth}p{.15\linewidth}p{.15\linewidth}}
    \# & \multicolumn{2}{c}{Player} && \multicolumn{2}{c}{Game Master} \\
    \hline

    \theutterance \stepcounter{utterance}  
    & & & \multicolumn{4}{p{0.6\linewidth}}{
        \cellcolor[rgb]{0.9,0.9,0.9}{
            \makecell[{{p{\linewidth}}}]{
                \texttt{\tiny{[P1$\langle$GM]}}
                \texttt{You are an expert in causal inference. The following question is not a typical commonsense query, but rather a meticulously designed question created by a professor specializing in causal inference, intended to assess the students' mastery of the course content.} \\
\texttt{Imagine a self{-}contained, hypothetical world with only the following conditions, and without any unmentioned factors or causal relationships: Gender has a direct effect on smoking and lung cancer. Smoking has a direct effect on tar deposit. Tar deposit has a direct effect on lung cancer. Gender is unobserved. For nonsmokers and with no tar deposit, the probability of lung cancer is 79\%. For nonsmokers and with high tar deposit, the probability of lung cancer is 49\%. For smokers and with no tar deposit, the probability of lung cancer is 71\%. For smokers and with high tar deposit, the probability of lung cancer is 40\%. For nonsmokers, the probability of high tar deposit is 79\%. For smokers, the probability of high tar deposit is 40\%. For smokers, would it be more likely to see lung cancer if the person had been a nonsmoker?} \\
\texttt{Start your answer with 'Yes' or 'No', followed by additional reasoning or evidence to support your explanation.} \\
            }
        }
    }
     \\ \\

    \theutterance \stepcounter{utterance}  
    & \multicolumn{4}{p{0.6\linewidth}}{
        \cellcolor[rgb]{0.561,0.737,0.561}{
            \makecell[{{p{\linewidth}}}]{
                \texttt{\tiny{[P1$\rangle$GM]}}
                \texttt{No} \\
\\ 
\texttt{To determine whether a smoker would be more likely to see lung cancer if they had been a nonsmoker, we must calculate the counterfactual probability of lung cancer for a smoker under the intervention of being a nonsmoker (\$P(Y=1 | do(S=0))\$) and compare it to the observed probability of lung cancer for a smoker (\$P(Y=1 | S=1)\$).} \\
\\ 
\texttt{**1. Define the Variables and Probabilities**} \\
\texttt{Let \$G\$ be Gender (unobserved), \$S\$ be Smoking (\$1=\$ smoker, \$0=\$ nonsmoker), \$T\$ be Tar Deposit (\$1=\$ high, \$0=\$ no), and \$Y\$ be Lung Cancer (\$1=\$ yes, \$0=\$ no).} \\
\\ 
\texttt{The problem provides the following conditional probabilities:} \\
\texttt{*   \$P(Y=1 | S=0, T=0) = 0.79\$} \\
\texttt{*   \$P(Y=1 | S=0, T=1) = 0.49\$} \\
\texttt{*   \$P(Y=1 | S=1, T=0) = 0.71\$} \\
\texttt{*   \$P(Y=1 | S=1, T=1) = 0.40\$} \\
\texttt{*   \$P(T=1 | S=0) = 0.79\$} \\
\texttt{*   \$P(T=1 | S=1) = 0.40\$} \\
\\ 
\texttt{**2. Calculate the Observed Probability for Smokers (\$P(Y=1 | S=1)\$)**} \\
\texttt{Using the law of total probability, we sum over the possible states of Tar (\$T\$):} \\
\texttt{\$\$P(Y=1 | S=1) = P(Y=1 | S=1, T=0)P(T=0 | S=1) + P(Y=1 | S=1, T=1)P(T=1 | S=1)\$\$} \\
\\ 
\texttt{Given \$P(T=1 | S=1) = 0.40\$, then \$P(T=0 | S=1) = 1 {-} 0.40 = 0.60\$.} \\
\texttt{\$\$P(Y=1 | S=1) = (0.71 \textbackslash{}times 0.60) + (0.40 \textbackslash{}times 0.40)\$\$} \\
\texttt{\$\$P(Y=1 | S=1) = 0.426 + 0.160 = 0.586\$\$} \\
\\ 
\texttt{**3. Calculate the Counterfactual Probability (\$P(Y=1 | do(S=0))\$)**} \\
\texttt{The question asks: "For smokers, would it be more likely to see lung cancer if the person had been a nonsmoker?" This requires calculating the probability of lung cancer if we intervene to set Smoking to 0 (\$do(S=0)\$).} \\
\\ 
\texttt{In this hypothetical world, the causal structure is \$G \textbackslash{}to S\$, \$G \textbackslash{}to Y\$, \$S \textbackslash{}to T\$, and \$T \textbackslash{}to Y\$.} \\
\texttt{When we perform the intervention \$do(S=0)\$, we break the arrow from \$G\$ to \$S\$. However, the distribution of the unobserved confounder \$G\$ remains the same as the background population.} \\
\texttt{Crucially, the probability of Tar given the intervention \$do(S=0)\$ is simply \$P(T | S=0)\$ because \$S\$ is the direct parent of \$T\$. The distribution of \$T\$ under the intervention is determined by the mechanism \$P(T|S)\$.} \\
\texttt{\$\$P(T=1 | do(S=0)) = P(T=1 | S=0) = 0.79\$\$} \\
\texttt{\$\$P(T=0 | do(S=0)) = 1 {-} 0.79 = 0.21\$\$} \\
\\ 
\texttt{Now we calculate the probability of Lung Cancer under this intervention:} \\
\texttt{\$\$P(Y=1 | do(S=0)) = P(Y=1 | S=0, T=0)P(T=0 | do(S=0)) + P(Y=1 | S=0, T=1)P(T=1 | do(S=0))\$\$} \\
\texttt{\$\$P(Y=1 | do(S=0)) = (0.79 \textbackslash{}times 0.21) + (0.49 \textbackslash{}times 0.79)\$\$} \\
\texttt{\$\$P(Y=1 | do(S=0)) = 0.1659 + 0.3871 = 0.553\$\$} \\
\\ 
\texttt{*Note: Even though Gender is unobserved, the calculation of the Average Treatment Effect (ATE) or the counterfactual probability for the population relies on the observed conditional probabilities of the mediators and outcomes, provided the Markov condition holds and we are calculating the marginal effect over the population distribution of \$G\$.*} \\
\\ 
\texttt{**4. Comparison**} \\
\texttt{*   Observed probability for smokers: **0.586**} \\
\texttt{*   Counterfactual probability (if they were nonsmokers): **0.553**} \\
\\ 
\texttt{The question asks: "For smokers, would it be more likely to see lung cancer if the person had been a nonsmoker?"} \\
\texttt{This translates to: Is \$P(Y=1 | do(S=0)) > P(Y=1 | S=1)\$?} \\
\texttt{\$0.553 > 0.586\$ is **False**.} \\
\\ 
\texttt{In this specific hypothetical scenario, the "nonsmoker" state is associated with a higher probability of tar deposits (79\% vs 40\%), and the baseline risk of lung cancer for nonsmokers with no tar is very high (79\%). The combination of these factors results in a higher overall risk for the "nonsmoker" counterfactual state compared to the "smoker" state in this specific dataset, despite the direct causal arrow from smoking to tar. The high baseline risk of lung cancer in the nonsmoking group (likely driven by the unobserved Gender variable) outweighs the protective effect of having less tar.} \\
\\ 
\texttt{Therefore, a smoker would **not** be more likely to see lung cancer if they had been a nonsmoker; they would actually be slightly less likely.} \\
            }
        }
    }
    & & \\ \\

    \theutterance \stepcounter{utterance}  
    & & & \multicolumn{2}{p{0.3\linewidth}}{
        \cellcolor[rgb]{0.9,0.9,0.9}{
            \makecell[{{p{\linewidth}}}]{
                \texttt{\tiny{[GM$|$GM]}}
                \texttt{No} \\
            }
        }
    }
    & & \\ \\

    \theutterance \stepcounter{utterance}  
    & & & \multicolumn{2}{p{0.3\linewidth}}{
        \cellcolor[rgb]{0.9,0.9,0.9}{
            \makecell[{{p{\linewidth}}}]{
                \texttt{\tiny{[GM$|$GM]}}
                \texttt{no} \\
            }
        }
    }
    & & \\ \\

    \theutterance \stepcounter{utterance}  
    & & & \multicolumn{2}{p{0.3\linewidth}}{
        \cellcolor[rgb]{0.9,0.9,0.9}{
            \makecell[{{p{\linewidth}}}]{
                \texttt{\tiny{[GM$|$GM]}}
                \texttt{game\_result = WIN} \\
            }
        }
    }
    & & \\ \\

\end{supertabular}
}

\end{document}
